\heading{理论力学-第2次作业}



\begin{question}
以 \(L = L(q, \dot q, \chi, t)\) 为例，证明把辅助变量 \(\chi = \chi(q, \dot q, t)\) 代入原 Lagrange 量得到的有效 Lagrange 量
\[
L_{\text{eff}}(q,\dot q,t) \equiv L\bigl(q,\dot q,\chi(q,\dot q,t),t\bigr)
\]
能给出 \(q\) 的运动方程。
\begin{solution}
原 Lagrange 量 \(L(q,\dot q,\chi,\dot\chi,t)\) 的变分给出两个 Euler--Lagrange 方程：
\begin{align*}
\frac{\mathrm d}{\mathrm dt}\frac{\partial L}{\partial\dot q} - \frac{\partial L}{\partial q} &= 0, \tag{1}\\
\frac{\mathrm d}{\mathrm dt}\frac{\partial L}{\partial\dot\chi} - \frac{\partial L}{\partial\chi} &= 0. \tag{2}
\end{align*}
将 \(\chi = \chi(q,\dot q,t)\) 代入后，有效 Lagrange 量对 \(q\) 的变分为
\[
\delta\int L_{\text{eff}}\,\mathrm dt
= \int\Bigl(
\frac{\partial L}{\partial q}\delta q
+ \frac{\partial L}{\partial\dot q}\delta\dot q
+ \frac{\partial L}{\partial\chi}\delta\chi
+ \frac{\partial L}{\partial\dot\chi}\delta\dot\chi
\Bigr)\mathrm dt,
\]
其中
\[
\delta\chi = \frac{\partial\chi}{\partial q}\delta q + \frac{\partial\chi}{\partial\dot q}\delta\dot q,\qquad
\delta\dot\chi = \frac{\mathrm d}{\mathrm dt}\delta\chi .
\]
对 \(\delta\dot q\) 和 \(\delta\dot\chi\) 项分部积分，并利用 \(\delta\chi\) 的表达式，得到
\[
\int\Bigl[
\frac{\partial L}{\partial q} - \frac{\mathrm d}{\mathrm dt}\frac{\partial L}{\partial\dot q}
+ \frac{\partial L}{\partial\chi}\frac{\partial\chi}{\partial q}
- \frac{\mathrm d}{\mathrm dt}\Bigl(\frac{\partial L}{\partial\dot\chi}\frac{\partial\chi}{\partial\dot q}\Bigr)
+ \frac{\partial L}{\partial\dot\chi}\frac{\partial^2\chi}{\partial q\partial\dot q}\dot q
+ \cdots
\Bigr]\delta q\,\mathrm dt = 0 .
\]
整理后，含 \(\chi\) 偏导的项为
\[
\frac{\partial L}{\partial\chi} - \frac{\mathrm d}{\mathrm dt}\frac{\partial L}{\partial\dot\chi}
\]
乘以 \(\chi\) 对 \(q\) 或 \(\dot q\) 的偏导再积分，而式 (2) 恰好使这些项为零。因此只剩下
\[
\frac{\mathrm d}{\mathrm dt}\frac{\partial L}{\partial\dot q} - \frac{\partial L}{\partial q} = 0,
\]
即由 \(L_{\text{eff}}\) 得到的 \(q\) 运动方程与原完整体系一致。
\end{solution}
\end{question}

\begin{question}
均匀护栏铁索不可拉伸，长为 \(L\)，质量为 \(m\)，两端等高，间距为 \(d\)。求：
\begin{itemize}
  \item[\textup{(i)}] 选择合适坐标，求重力势能 \(V\) 作为形状的泛函；
  \item[\textup{(ii)}] 用 Lagrange 乘子法求形状满足的 Euler--Lagrange 方程；
  \item[\textup{(iii)}] 求解 E--L 方程，证明解是悬链线。
\end{itemize}

\pyfig[0.55\linewidth]{figures/hw2_fig1.png}{两端固定的均匀铁索（悬链线）示意图}
\begin{solution}
取坐标系原点在左端点水平位置，\(y\) 轴向下。设铁索形状为 \(y(x)\)，两端条件 \(y(0)=y(d)=0\)。

\textbf{(i)} 线密度 \(\lambda = m/L\)。弧长微元 \(\mathrm ds = \sqrt{1+y'^2}\,\mathrm dx\)，总长约束为
\[
\int_0^d \sqrt{1+y'^2}\,\mathrm dx = L .
\]
以 \(y=0\) 为零势点，重力势能泛函为
\[
V[y] = -\lambda g \int_0^d y\,\mathrm ds
= -\frac{mg}{L}\int_0^d y\sqrt{1+y'^2}\,\mathrm dx .
\]

\textbf{(ii)} 引入 Lagrange 乘子 \(\mu\)（常数），构造
\[
F = -\lambda g\, y\sqrt{1+y'^2} + \mu\sqrt{1+y'^2}
= \bigl(-\lambda g\, y + \mu\bigr)\sqrt{1+y'^2}.
\]
Euler--Lagrange 方程
\[
\frac{\mathrm d}{\mathrm dx}\frac{\partial F}{\partial y'} - \frac{\partial F}{\partial y} = 0
\]
给出
\[
\frac{\partial F}{\partial y} = -\lambda g\sqrt{1+y'^2},
\qquad
\frac{\partial F}{\partial y'} = \frac{(-\lambda g\, y + \mu)\,y'}{\sqrt{1+y'^2}} .
\]
代入得
\[
\frac{\mathrm d}{\mathrm dx}\Bigl[\frac{(-\lambda g\, y + \mu)\,y'}{\sqrt{1+y'^2}}\Bigr]
+ \lambda g\sqrt{1+y'^2} = 0 .
\]

\textbf{(iii)} 上式可写为全微分形式。注意到
\[
\frac{\mathrm d}{\mathrm dx}\Bigl[\frac{-\lambda g\, y + \mu}{\sqrt{1+y'^2}}\Bigr]
= \frac{(-\lambda g\, y')\sqrt{1+y'^2}
- (-\lambda g\, y + \mu)\frac{y'y''}{\sqrt{1+y'^2}}}{1+y'^2}
+ \lambda g\sqrt{1+y'^2}\cdot\text{(?)}
\]
更直接地，由 Euler--Lagrange 方程一次积分：
令 \(F\) 不显含 \(x\)，故 Beltrami 恒等式给出
\[
F - y'\frac{\partial F}{\partial y'} = \text{常数}.
\]
计算
\[
F - y'\frac{\partial F}{\partial y'}
= (-\lambda g\, y + \mu)\sqrt{1+y'^2}
- \frac{(-\lambda g\, y + \mu)\,y'^2}{\sqrt{1+y'^2}}
= \frac{-\lambda g\, y + \mu}{\sqrt{1+y'^2}} = C,
\]
其中 \(C\) 为积分常数。于是
\[
\sqrt{1+y'^2} = \frac{-\lambda g\, y + \mu}{C}
\quad\Longrightarrow\quad
y'^2 = \Bigl(\frac{-\lambda g\, y + \mu}{C}\Bigr)^2 - 1 .
\]
令 \(u = -\lambda g\, y + \mu\)，则
\[
\frac{\mathrm du}{\mathrm dx} = -\lambda g\, y'
= -\lambda g\,\sqrt{\bigl(u/C\bigr)^2 - 1},
\qquad
\frac{\mathrm du}{\sqrt{u^2/C^2 - 1}} = -\lambda g\,\mathrm dx .
\]
积分得
\[
C\cosh^{-1}\!\Bigl(\frac{u}{C}\Bigr) = -\lambda g\, x + C_1,
\]
即
\[
u = C\cosh\!\Bigl(\frac{C_1 - \lambda g\, x}{C}\Bigr).
\]
代回 \(y\)：
\[
y(x) = \frac{C}{\lambda g}\Bigl[\cosh\!\Bigl(\frac{C_1}{C}\Bigr)
- \cosh\!\Bigl(\frac{C_1 - \lambda g\, x}{C}\Bigr)\Bigr],
\]
其中常数由边界条件 \(y(0)=y(d)=0\) 和弧长约束 \(\int_0^d \sqrt{1+y'^2}\,\mathrm dx = L\) 确定。上式即为标准悬链线（catenary）形式：
\[
\boxed{\;y(x) = a\Bigl[\cosh\!\Bigl(\frac{x_0}{a}\Bigr)
- \cosh\!\Bigl(\frac{x-x_0}{a}\Bigr)\Bigr]\;},
\]
其中 \(a = C/(\lambda g)\)，且常数由边界和总长决定。
\end{solution}
\end{question}

\begin{question}
一半径为 \(a\) 的光滑圆环，绕平面内过圆心的铅直轴以恒定的角速度 \(\omega\) 转动，质量为 \(m\) 的粒子沿环运动。取 \(\theta\) 为广义坐标（\(\theta\) 为粒子与竖直方向的夹角）。

\pyfig[0.4\linewidth]{figures/hw2_q3_diagram.png}{绕铅直轴匀速转动的光滑圆环}
\begin{itemize}
  \item[\textup{(i)}] 写出粒子的 Lagrange 量；
  \item[\textup{(ii)}] 证明粒子的 Hamiltonian 为运动常数，而总能量 \(E\) 不是，并讨论其原因。
\end{itemize}
\begin{solution}
取圆环中心为原点，竖直轴为 \(z\) 轴。圆环绕 \(z\) 轴以角速度 \(\omega\) 匀速转动，粒子沿环运动。粒子在空间中的坐标为
\[
x = a\sin\theta\cos(\omega t),\qquad
y = a\sin\theta\sin(\omega t),\qquad
z = -a\cos\theta .
\]
（取 \(z=0\) 为环心高度，\(z\) 向上为正。）

\textbf{(i)} 速度分量为
\begin{align*}
\dot x &= a\dot\theta\cos\theta\cos(\omega t) - a\omega\sin\theta\sin(\omega t),\\
\dot y &= a\dot\theta\cos\theta\sin(\omega t) + a\omega\sin\theta\cos(\omega t),\\
\dot z &= a\dot\theta\sin\theta .
\end{align*}
动能为
\[
T = \frac12 m(\dot x^2 + \dot y^2 + \dot z^2)
= \frac12 m\Bigl(a^2\dot\theta^2 + a^2\omega^2\sin^2\theta\Bigr).
\]
势能以环心为零点：
\[
V = mgz = -mg a\cos\theta .
\]
故 Lagrange 量为
\[
\boxed{\;L = T - V
= \frac12 m a^2\dot\theta^2 + \frac12 m a^2\omega^2\sin^2\theta + mg a\cos\theta\;}.
\]

\textbf{(ii)} 广义动量
\[
p_\theta = \frac{\partial L}{\partial\dot\theta} = m a^2\dot\theta .
\]
Hamiltonian 为
\[
H = p_\theta\dot\theta - L
= m a^2\dot\theta^2
- \Bigl(\frac12 m a^2\dot\theta^2 + \frac12 m a^2\omega^2\sin^2\theta + mg a\cos\theta\Bigr)
= \frac12 m a^2\dot\theta^2 - \frac12 m a^2\omega^2\sin^2\theta - mg a\cos\theta .
\]
由于 \(L\) 不显含时间 \(t\)，Hamiltonian 是运动常数（守恒量）：
\[
\boxed{\;\frac{\mathrm d H}{\mathrm d t} = 0\;}.
\]
总能量为
\[
E = T + V
= \frac12 m a^2\dot\theta^2 + \frac12 m a^2\omega^2\sin^2\theta - mg a\cos\theta .
\]
比较可知
\[
E = H + m a^2\omega^2\sin^2\theta .
\]
由于 \(\theta\) 随时间变化，\(E\) 不是常数。其原因在于：圆环的转动对粒子做功，系统存在非保守力作用，机械能不守恒。旋转环提供的约束是非定常约束，广义力包含与转动相关的惯性力（离心力和 Coriolis 力），因此总能量 \(E\) 不守恒；而 Hamiltonian（在 \(L\) 不显含时间时）始终是运动常数。
\end{solution}
\end{question}

\begin{question}
一维阻尼振子的 Lagrange 量为
\[
L = e^{\lambda t}\Bigl(\frac12 m\dot q^2 - \frac12 m\omega^2 q^2\Bigr),
\]
其中 \(m,\omega,\lambda\) 是常数。
\begin{itemize}
  \item[\textup{(i)}] 取常数 \(\alpha\)，证明在变换
  \[
  t \to \tilde t = t + \alpha,\qquad
  q(t) \to \tilde q(\tilde t) = e^{-\lambda\alpha/2}\, q(t)
  \]
  下，\(\tilde S[\tilde q] = S[q]\)；
  \item[\textup{(ii)}] 考虑 \(\alpha\) 为无穷小量，用 Noether 定理求该对称变换对应的运动常数。
\end{itemize}
\begin{solution}
\textbf{(i)} 作用量为
\[
S[q] = \int L\,\mathrm dt
= \int e^{\lambda t}\Bigl(\frac12 m\dot q^2 - \frac12 m\omega^2 q^2\Bigr)\mathrm dt .
\]
变换后的作用量为
\[
\tilde S[\tilde q] = \int e^{\lambda\tilde t}
\Bigl(\frac12 m\frac{\mathrm d\tilde q}{\mathrm d\tilde t}^2
- \frac12 m\omega^2 \tilde q^2\Bigr)\mathrm d\tilde t .
\]
由于 \(\tilde t = t + \alpha\)，有 \(\mathrm d\tilde t = \mathrm dt\)，且 \(\tilde q(\tilde t) = e^{-\lambda\alpha/2} q(t)\)。因此
\[
\frac{\mathrm d\tilde q}{\mathrm d\tilde t}
= e^{-\lambda\alpha/2}\frac{\mathrm d q}{\mathrm d t},
\qquad
\tilde q = e^{-\lambda\alpha/2} q(t) .
\]
代入：
\begin{align*}
\tilde S[\tilde q]
&= \int e^{\lambda(t+\alpha)}\Bigl(
\frac12 m\, e^{-\lambda\alpha}\dot q^2
- \frac12 m\omega^2 e^{-\lambda\alpha} q^2\Bigr)\mathrm dt \\
&= \int e^{\lambda t}\Bigl(\frac12 m\dot q^2 - \frac12 m\omega^2 q^2\Bigr)\mathrm dt
= S[q] .
\end{align*}
故作用量在该变换下不变，该变换是 Noether 对称变换。

\textbf{(ii)} 无穷小变换。令 \(\alpha \to 0\)，则
\[
\delta t = \tilde t - t = \alpha,\qquad
\delta q = \tilde q(\tilde t) - q(t) = e^{-\lambda\alpha/2} q(t) - q(t)
\approx -\frac{\lambda\alpha}{2} q + \dot q\,\alpha .
\]
Noether 定理：对称变换 \((\delta t,\delta q)\) 对应的守恒量为
\[
I = \frac{\partial L}{\partial\dot q}\bigl(\delta q - \dot q\,\delta t\bigr) + L\,\delta t .
\]
计算
\[
\frac{\partial L}{\partial\dot q} = e^{\lambda t}\, m\dot q .
\]
代入：
\begin{align*}
I &= e^{\lambda t} m\dot q\Bigl(-\frac{\lambda\alpha}{2} q + \dot q\alpha - \dot q\alpha\Bigr)
+ \alpha\, e^{\lambda t}\Bigl(\frac12 m\dot q^2 - \frac12 m\omega^2 q^2\Bigr) \\
&= e^{\lambda t} m\dot q\Bigl(-\frac{\lambda\alpha}{2} q\Bigr)
+ \alpha\, e^{\lambda t}\Bigl(\frac12 m\dot q^2 - \frac12 m\omega^2 q^2\Bigr) \\
&= \alpha\, e^{\lambda t}\Bigl(\frac12 m\dot q^2 - \frac12 m\omega^2 q^2
- \frac{\lambda m}{2} q\dot q\Bigr).
\end{align*}
去掉常数因子 \(\alpha\)，Noether 守恒量为
\[
\boxed{\;I = e^{\lambda t}\Bigl(\frac12 m\dot q^2 - \frac12 m\omega^2 q^2
- \frac{\lambda m}{2} q\dot q\Bigr)\;}.
\]
可以验证 \(\mathrm d I/\mathrm d t = 0\) 正是阻尼振子的运动方程。
\end{solution}
\end{question}

\begin{question}
证明变换
\[
t \to \tilde t = \frac{\alpha t + \beta}{\gamma t + \delta},\qquad
q(t) \to \tilde q(\tilde t) = \frac{q(t)}{\gamma t + \delta},
\qquad
\det\begin{pmatrix}\alpha & \beta \\ \gamma & \delta\end{pmatrix} = 1
\]
是作用量
\[
S = \int \Bigl(\frac12 m\dot q^2 - \frac{\lambda}{q^2}\Bigr)\mathrm dt
\]
的对称变换；其中 \(\alpha,\beta,\gamma,\delta,m,\lambda\) 是常数。
\begin{solution}
需证明 \(\tilde S[\tilde q] = S[q]\)。变换前后时间坐标的关系为
\[
\tilde t = \frac{\alpha t + \beta}{\gamma t + \delta},
\qquad
\mathrm d\tilde t = \frac{\alpha\delta - \beta\gamma}{(\gamma t + \delta)^2}\,\mathrm dt
= \frac{\mathrm dt}{(\gamma t + \delta)^2},
\]
其中利用了 \(\det = \alpha\delta - \beta\gamma = 1\)。

定义 \(Q(t) = q(t)\)，则 \(\tilde q(\tilde t) = Q(t)/(\gamma t + \delta)\)。计算 \(\tilde q\) 对 \(\tilde t\) 的导数：
\[
\frac{\mathrm d\tilde q}{\mathrm d\tilde t}
= \frac{\mathrm d}{\mathrm d\tilde t}\!\left(\frac{Q}{\gamma t + \delta}\right)
= \frac{\mathrm d t}{\mathrm d\tilde t}\,
\frac{\mathrm d}{\mathrm d t}\!\left(\frac{Q}{\gamma t + \delta}\right)
= (\gamma t + \delta)^2\,
\frac{\dot Q(\gamma t + \delta) - \gamma Q}{(\gamma t + \delta)^2}
= \dot Q(\gamma t + \delta) - \gamma Q .
\]

变换后的作用量为
\begin{align*}
\tilde S[\tilde q]
&= \int\Bigl(\frac12 m\Bigl(\frac{\mathrm d\tilde q}{\mathrm d\tilde t}\Bigr)^2
- \frac{\lambda}{\tilde q^2}\Bigr)\mathrm d\tilde t \\
&= \int\Bigl[
\frac12 m\bigl(\dot Q(\gamma t + \delta) - \gamma Q\bigr)^2
- \lambda (\gamma t + \delta)^2 \frac{1}{Q^2}
\Bigr]\frac{\mathrm dt}{(\gamma t + \delta)^2} .
\end{align*}
展开第一项：
\[
\bigl(\dot Q(\gamma t + \delta) - \gamma Q\bigr)^2
= \dot Q^2(\gamma t + \delta)^2 - 2\gamma Q\dot Q(\gamma t + \delta) + \gamma^2 Q^2 .
\]
代入积分：
\[
\tilde S[\tilde q]
= \int\Bigl[
\frac12 m\dot Q^2
- m\gamma\,\frac{Q\dot Q}{\gamma t + \delta}
+ \frac12 m\gamma^2\,\frac{Q^2}{(\gamma t + \delta)^2}
- \frac{\lambda}{Q^2}
\Bigr]\mathrm dt .
\]
注意到
\[
- m\gamma\,\frac{Q\dot Q}{\gamma t + \delta}
+ \frac12 m\gamma^2\,\frac{Q^2}{(\gamma t + \delta)^2}
\]
为全微分项。具体地，
\[
\frac{\mathrm d}{\mathrm dt}\!\left(-\frac12 m\gamma\,\frac{Q^2}{\gamma t + \delta}\right)
= -\frac12 m\gamma\,\frac{2Q\dot Q(\gamma t + \delta) - \gamma Q^2}{(\gamma t + \delta)^2}
= - m\gamma\,\frac{Q\dot Q}{\gamma t + \delta}
+ \frac12 m\gamma^2\,\frac{Q^2}{(\gamma t + \delta)^2} .
\]
因此
\[
\tilde S[\tilde q]
= \int\Bigl(\frac12 m\dot Q^2 - \frac{\lambda}{Q^2}\Bigr)\mathrm dt
+ \Bigl[-\frac12 m\gamma\,\frac{Q^2}{\gamma t + \delta}\Bigr]_{\text{端点}} .
\]
在端点固定的变分原理中，边界项对运动方程无贡献（在端点处变分为零，边界项不影响变分极值）。忽略边界项后得
\[
\tilde S[\tilde q] = \int\Bigl(\frac12 m\dot q^2 - \frac{\lambda}{q^2}\Bigr)\mathrm dt = S[q],
\]
故该变换为作用量的对称变换。
\end{solution}
\end{question}
